\chapter{Gastric Lavage}

\section*{Overview}
Gastric lavage is the washing out of the stomach using a large-bore tube with instillation and aspiration of fluid. It is used to remove ingested toxins or to manage gastrointestinal bleeding.

\section*{Alternative Names}
Stomach pumping, Gastric irrigation

\section*{Principle}
A wide-bore orogastric tube is passed into the stomach, and small aliquots of fluid, typically warm saline, are instilled and aspirated in repeated cycles. This mechanically removes gastric contents and unabsorbed toxin. Removal depends on mechanical clearance rather than neutralization or absorption of the substance.

\section*{History}
Gastric lavage has been performed for more than a century as a method of gastric decontamination after poisoning.

\section*{Why It Is Done}
Gastric lavage is ordered for potentially life-threatening ingestions when performed soon after the event, ideally within one hour, and when the substance is not contraindicated. It may also be used therapeutically for gastrointestinal bleeding or to obtain gastric contents for analysis.

\section*{Is This a Primary or Secondary Test or Procedure}
Secondary procedure of declining use; it is reserved for selected toxic ingestions rather than routine poisoning management.

\section*{How It Is Performed}
The patient is positioned head down on the left side, and a lubricated large-bore orogastric tube is passed orally into the stomach. After confirming tube position, 200--300\,mL aliquots of warm saline are instilled and aspirated until the returns are clear. The tube is then removed.

\section*{Preparation}
Airway protection is essential; patients with impaired consciousness or a depressed gag reflex are intubated first. The substance, time of ingestion, and vital signs are documented, and informed consent is obtained when possible.

\section*{Contraindications and Precautions}
Avoid lavage after ingestion of hydrocarbons, corrosives, or substances with high aspiration risk, and avoid sharp objects. It is contraindicated in esophageal varices, recent esophageal surgery, and severe bleeding diatheses.

\section*{Risks}
Aspiration pneumonia, laryngospasm, and esophageal or gastric perforation. Epistaxis and electrolyte disturbances from fluid shifts may also occur.

\section*{Aftercare and Recovery}
The patient is observed for complications such as aspiration, and supportive care continues. Poisoning management proceeds with activated charcoal, antidotes, and laboratory monitoring as indicated.

\section*{Outcome and Follow-up}
Lavage serves decontamination rather than diagnosis; effectiveness depends on timing, quantity, and type of toxin. Clinical assessment and toxicologic monitoring determine the outcome.

\section*{Expected Results}
Clear return of lavage fluid with no further recovery of gastric contents; no quantitative diagnostic value is expected.

\section*{Follow-up}
Observation for delayed toxicity with supportive laboratory monitoring and serial clinical assessment. Referral to toxicology services and psychiatric evaluation is appropriate after intentional ingestion.

\section*{References}
\begin{enumerate}
  \item MedlinePlus. Gastric lavage. U.S. National Library of Medicine; 2025.
  \item Current Medical Diagnosis and Treatment. McGraw-Hill; 2025.
\end{enumerate}

\clearpage
